\documentclass{fmfkpb}

\avtor{Luka Lavš}
\naslov{Mooreovi grafi in posplošeni večkotniki}
\letnica{2026}
\mentor{doc. dr. Janoš Vidali}

\newcommand{\red}{\operatorname{red}}
\newcommand{\srg}{\operatorname{srg}}
\newcommand{\diam}{diam}
\newcommand{\N}{\mathbb{N}}
\newcommand{\ZZ}{\mathbb{Z}}
\newcommand{\reach}{\operatorname{Reach}}
\newcommand{\gr}{\mathscr{G}}

\newcommand{\A}{\mathcal{A}}
\newcommand{\II}{\mathbf{I}}
\newcommand{\IIw}[2]{#1 \;\II\; #2}
\newcommand{\IIwd}[2]{#1 \;\II^D\; #2}
\newcommand{\PP}{\mathcal{P}}
\newcommand{\LL}{\mathcal{L}}

\begin{document}

\makefrontpage
\cleardoublepage

\section{Mooreovi grafi in posplošeni večkotniki}

\paragraph{Par najosnovnejših pojmov}
V teoriji grafov, rečemo, da je $D$ premer grafa, če je največja razdalja med katerima
koli vozliščema enaka $D$.
Dalje pravimo, da ima graf ožino $g$, če je to dolžina najkrajšega cikla.
Grafu, v katerem so stopnje vseh vozlišč enake $\Delta$, pravimo $\Delta$-regularen.


\paragraph{Mooreova meja in Mooreov graf}
Definiramo zdaj Mooreovo mejo kot maksimalno število vozlišč, ki bi jih 
graf z danim premerom in dano maksimalno stopnjo, 
v optimističnem primeru, lahko imel.

Vidimo, da se
vsak graf začne v nekem vozlišču ($1 + \ldots$),
ki ima največ $\Delta$ sosedov ($1 + \Delta + \ldots$), 
in vsak izmed teh sosedov
pa ima največ $\Delta - 1$ novih sosedov ($1 + \Delta + \Delta(\Delta - 1) + \ldots$).
Vidimo, da lahko tako nadaljujemo dokler nismo primorani vejanja končati 
zaradi premera $D$. Dobili smo 
($1 + \Delta + \Delta(\Delta - 1) + \ldots \Delta (\Delta - 1)^{D - 1}$) in 
temu rečemo Mooreova meja.

Tu je morda smiselno omeniti, da smo pri konstrukciji Mooreove meje ustvarili 
drevesni podgraf, ki je nekakšen "idealni" podgraf, saj ima največje 
možno število vozlišč za dan premer in maksimalno stopnjo.
Ni pa nujno, da za vsak premer in maksimalno stopnjo, tak graf tudi obstaja,
vendar če, mu rečemo Mooreov graf.
V resnici so taki grafi zelo redki.

\paragraph{Trivialni Mooreovi grafi}
Zdaj bom navedel vse trivialne primere Mooreovih grafov. Morda lahko 
začnem z najenostavnejšim, in sicer z trikotnikom. Trikotnik je Mooreov graf
za dan premer $1$ in maksimalno stopnjo $2$. (Narišem trikotnik in utemeljim,
da doseže Mooreovo mejo $1 + \Delta = 3$).

Dodatno lahko hitro ugotovimo, da so vsi cikli $C_{2D + 1}$ Mooreovi grafi 
za poljuben premer $D$ in maksimalno stopnjo $2$.
Preostali trivialni Mooreovi grafi pa so še polni grafi $K_{\Delta + 1}$, 
za poljubno stopnjo in premer $1$. (Narišem polni graf kot polno povezan graf).

\paragraph{Lastnosti Mooreovih grafov}
Kot sem že namignil, so ostali netrivialni Mooreovi grafi zelo redki, 
zato je smiselno analizirati njihove lastnosti, da bi razumeli zakaj je temu
tako, in da bi jih morda celo našli.

Z enostavnejšima premislekoma se izkaže, da so Mooreovi grafi regularni,
in da je dolžina njihovega najkrajšega cikla, in torej ožina, enaka $g=2D+1$.
Posledično tudi obstaja natanko ena pot dolžine največ $D$, med vsakima vozliščema.

Opremljeni s tem dejstvom lahko potem tvorimo matrično enačbo, kateri 
morajo vsi Mooreovi grafi zadoščati. In sicer je ta enačba navedena na prosojnici,
kjer je morda vredno dodati, da je tisti $G$ matrika katere 
vhodi merijo število poti dolžine največ $D$ med vozlišči, tisti "J" pa 
je matrika samih enic ustrezne dimenzije.

\paragraph{Mooreov graf premera 2}
Mimogrede, Mooreov graf premera $2$ je tudi krepko regularen, sosednja vozlišča imajo 
natanko nič skupnih sosedov, nesosednja pa natanko enega.

Ravno prej omenjena matrična enačba se v primeru ko je premer enak $2$ znatno
poenostavi, in sicer dobimo ($A^2 + A + I = J$), kjer je $A$ matrika 
sosednosti grafa in $J$ zopet matrika samih enic.

Ker iz linearne algebre vemo, da ima matrika samih enic lastni vrednosti
$\Delta^2 + 1$ in $0$, dobimo teli dve enačbi. (Pokažem na prosojnici).

Ko ju rešimo, dobimo možne lastne vrednosti matrike sosednosti $A$.
Zanimajo pa nas še njihove večkratnosti, in v ta namen uporabimo
še splošno znano enačbo ($0 = sl(A) = \sum\lambda_i$).

Analiziramo jo s pomočjo celosti, lihosti in izreka o ničlah 
celoštevilskih polinomov in dobimo, da so večkratnosti take (pokažem na prosojnici), 
in da so možne stopnje Mooreovih grafov premera $2$ zgolj $\{2, 3, 7, 57\}$.

\paragraph{Petersenov graf}
Morda nas bi zdaj zamikalo, da bi take grafe poskušali tudi konstruirati.
Za začetek se preizkusimo s tistim, ki ima stopnjo $3$.

Obstaja več različnih konstrukcij, danes bom pokazal dve.

Lahko ga konstruiramo iz drevesnega podgrafa, kjer poskušamo povezati zadnji
nivo vozlišč brez da bi kršili pogoj o ožini. 
(Narišem drevesni podgraf in povežem sosede).

Alternativno ga bi lahko konstruirali tudi iz cikla $C_5$, 
ki seveda zaradi pogoja o ožini obstaja. 
(Narišem pentagon in dokončam graf).

Mimogrede, ta graf se imenuje Petersenov graf.

\paragraph{Hoffman-Singletonov graf}
To ni bilo težko, poizkusimo zdaj konstruirati še naslednjega, tistega z stopnjo $7$.
Tukaj se graf konstruira ravno tako iz petkotnika, vendar pa je tak graf 
zdaj na $7\cdot 7 + 1 = 50$ vozliščih in precej kompleksnejši.

Na tej predstavitvi bom začrtal zgolj skico njegove konstrukcije, 
ki je tudi precej zahtevna. Je pa taka konstrukcija podrobno opisana v
moji diplomski nalogi in predstavlja tudi eno izmed daljših sekcij.

(Narišem pentagon, definiram nivoje $N(Q), N_2(Q)$ in gruče zgolj na hitro).
Izkaže se, da so si gruče paroma disjunktne; da če je en element
gruče povezan z elementom druge gruče, potem noben drug element prvotne 
gruče ne more biti povezan z katerim koli elementom druge gruče;
in da poljubna gruča ni povezana z gručama, ki sta sosednji glede na 
predstavnika. (Tu pokažem, da so $v_i \in Q$ taki predstavniki).

Posledično se izkaže tudi, da je risba, ki sem jo narisal do zdaj pravilna,
in da je stopnja v podgrafu induciranim z vozlišči prvega nivoja največ $2$.


Dodatno pa, da so vozlišča iz drugega nivoja povezana z natanko $5$ 
vozlišči iz prvega nivoja; vozlišča iz prvega nivoja pa natanko $4$-mi.
Tako se izkaže, da sta podgrafa inducirana z vozlišči prvega in 
drugega nivoja $2$-regularna, ter da sta oba podgrafa ravno
disjunktna unija ciklov dolžine večkratnika $5$.

Nekaj lem kasneje nam končno uspe pokazati, da je podgraf induciran z 
vozlišči prvega nivoja natanko unija $5$-ih $5$-ciklov, in da je
podgraf induciran z vozlišči drugega nivoja natanko unija $4$-ih $5$-ciklov.

Skupno ima tako iskani graf $5$ $5$-ciklov ene sorte in $5$ $5$-ciklov druge sorte.
Problem se poenostavi v iskanje $25$ permutacij, ki narekujejo povezave med 
takimi cikli.

Podrobnosti bom tu izpustil, lahko pa kot zanimivost omenim, da se na koncu
izkaže, da je graf generiran z $5$-mi cikli sorte $P$ in 
$5$-imi cikli sorte $Q$, kateri so med seboj 
povezani po pravilu $P^i_j \sim Q^k_{j + ik \pmod 5}$.

Za konec lahko še omenim, da se z uporabo Lagrangevega izreka o kosetih, 
ter izreka o orbitah in stabilizatorjih, ki sta ob izreka o grupah,
in še z nekaj truda, lahko določi število avtomorfizmov tega grafa,
ki je enako $50\cdot 7!$.
(Tu lahko omenim kaj izreka pravita, vendar najverjetneje ne bom.)

\paragraph{Graf valence 57}
Že tak graf je bilo kar zahtevno konstruirati, kako šele je pri tem grafu 
stopnje $57$. Tak graf bi moral stati na $1 +57\cdot 57$ vozliščih, kar je
neproporcionalno več od prejšjega. In tudi nikomur ga ni uspelo konstruirati,
niti zavrniti njegovega obstoja. Ne ve se če obstaja.

Se pa ve, da če obstaja, potem ni vozliščno tranzitiven, kar je precej nenavadno, 
glede na to da so vsi prejšnji grafi bili.
To se pokaže s pomočjo teorije grup, ki pokaže, da grupa avtomorfizmov 
vozliščno tranzitivnega grafa, takih parametrov, ne obstaja.

\paragraph{Obstoj Mooreovih grafov}
Mooreovi grafi, ki jih nisem navedel, ne obstajajo. Zato je smiselno problem 
razširiti in zajeti še tiste grafe, ki so sode ožine. Namreč dozdajšnja definicija
je dopuščala zgolj lihe ožine.

Definiramo razširjeno Mooreovo mejo kot spodnjo mejo števila vozlišč, ki bi jih
graf sode ožine moral imeti. (Narišem izpeljavo take meje, prek dvodelnosti.)

Taki grafi imajo zdaj sodo ožino, so tudi regularni in imajo določen premer, ki je 
enak polovici ožine. Takim grafom pravimo ravno incidenčni grafi simetričnih 
posplošenih večkotnikov.

\paragraph{Posplošeni večkotniki}
Posplošeni $n$-kotnik je neka geometrijska struktura, 
kjer imamo dvoje sort "vozlišč" -- točke in premice, 
ki so med seboj povezane z neko relacijo, tako, da cela struktura nima 
podgeometrijskih struktur, ki bi bile običajni manjši kotniki,
ter da se vsaka dva elementa nahajata v nekem običajnem $n$-kotniku.

Pravimo, da je posplošen $n$-kotnik simetričen, oziroma, da je reda $(s, s)$,
če je vsaka premica incidentna s $s + 1$ točkami, 
in vsaka točka incidentna s $s + 1$ premicami.

Obstaja trditev, ki tudi za splošnejši red izračuna število točk in premic
posplošenega $n$-kotnika, v našem simetričnem primeru pa vidimo, da se 
izračun tega števila natanko ujema z mejo, ki smo jo že navedli. Kakor bi 
tudi pričakovali.

Za primer lahko vzamemo najpreprostejši posplošen $n$-kotnik reda $(1, 1)$, in sicer 
je to kar običajni $n$-kotnik (narišem običajni $n$-kotnik, v smislu geometrije).
Posledično je vsak $2n$-kotnik posplošen Mooreov graf.

Izkaže se, da posplošeni $n$-kotniki reda večjega od $1$ obstajajo zgolj za $n \in \{2, 3, 4, 6, 8, 12\}$.
Primer ko je $n = 2$ je seveda posplošen dvokotnik, in njegov incidenčni graf
je polni dvodelni graf $K_{s, s}$. 


\paragraph{Projektivna ravnina}
Za zaključek lahko prikažem še zgled neke lepše projektivne ravnine, ki je 
posplošen simetrični trikotnik.
Projektivna ravnina je preprosto generirana nad nekim poljem $\mathbb{K}$, 
tako, da opazujemo vektorski prostor $V = \mathbb{K}^3 \setminus 0$, 
kjer točke definiramo kot vektorje, in premice kot ravnine. Relacijo sosednosti
pa definiramo glede na vsebovanost.

Tako lahko recimo za polje vzamemo končno polje $\mathbb{F}_2 = \{0, 1\}$,
Vektorski prostor brez ničle ima tako $2^3 - 1 = 7$ točk ter $7$ premic.
Točka leži na ravnini, če je produkt normale in točke enak $0$.
Tako recimo vidimo, da na ravnini, predstavljeni z normalo $(1, 0, 0)$ 
leži
vsaka točka, ki zadošča enačbi $1\cdot x + 0\cdot y + 0\cdot z = 0$. 
Torej vsaka izmed teh točk $\{(0,1,0),(0,0,1),(0,1,1)\}$ je sosedna z $(1, 0, 0)$.
Če bi to naredili še za vse ostale točke, bi dobili celoten posplošen 
simetrični trikotnik, ki mu pravimo Fanova ravnina.

\paragraph{Opomba sebi}
O ostalih večkotnikih bom razpravljal več v svoji diplomski nalogi,
s tem sem svojo dolgo predstavitev zaključil. 
(Lahko bi predstavil še zgled posplošenega $4$-kotnika, vendar ga nameravam bolj bežno omeniti,
saj bo drugače to verjetno že malo prekompleksno za omejen čas). Pokazal bom le še slike,
ki jih imam na prosojnicah.


\end{document}